\documentclass[11pt]{ctexart}
\usepackage[a4paper,top=2cm,bottom=2cm,left=2.2cm,right=2.2cm]{geometry}
\usepackage{fancyvrb}
\usepackage{xcolor}
\usepackage{colortbl}
\usepackage{fontspec}
\usepackage[most]{tcolorbox}
\usepackage{booktabs}
\usepackage{enumitem}
\usepackage{titlesec}
\usepackage{hyperref}

\setmonofont{Consolas}
\setCJKmonofont{SimHei}
\setlength{\parindent}{0pt}
\setlist{nosep,leftmargin=1.4em}

% ── 嵌入式（二）主题色：深石墨钢（工具链/机械感，与第一册同族区分） ──
\definecolor{deep}{RGB}{60,62,70}        % 主色：深石墨
\definecolor{accent}{RGB}{108,112,130}   % 强调：钢灰
\definecolor{keyline}{RGB}{90,130,200}   % 知识点：蓝
\definecolor{keybg}{RGB}{238,244,255}
\definecolor{warnline}{RGB}{200,120,60}  % 警告：橙
\definecolor{warnbg}{RGB}{255,243,235}
\definecolor{okline}{RGB}{26,116,88}     % 成果：绿
\definecolor{okbg}{RGB}{240,252,244}
\definecolor{tipline}{RGB}{132,136,152}  % 提示：浅钢灰
\definecolor{tipbg}{RGB}{242,243,246}
\definecolor{codebg}{RGB}{249,249,252}   % 代码底
\definecolor{rowtint}{RGB}{238,239,243}  % 表头浅钢灰

\titleformat{\section}{\large\bfseries\color{deep}}{\thesection}{0.6em}{}
\titleformat{\subsection}{\normalsize\bfseries\color{accent}}{\thesubsection}{0.5em}{}
\titleformat{\part}[display]{\centering\Huge\bfseries\color{deep}}{}{0pt}{\Huge}
\titlespacing{\section}{0pt}{14pt}{6pt}
\titlespacing{\subsection}{0pt}{10pt}{4pt}

\newtcolorbox{keybox}[1]{enhanced, colback=keybg, colframe=keyline,
  boxrule=0.6pt, arc=1.5mm, left=3mm, right=3mm, top=2mm, bottom=2mm,
  title={\bfseries #1}, fonttitle=\small\bfseries, coltitle=white,
  colbacktitle=keyline}
\newtcolorbox{warnbox}[1]{enhanced, colback=warnbg, colframe=warnline,
  boxrule=0.6pt, arc=1.5mm, left=3mm, right=3mm, top=2mm, bottom=2mm,
  title={\bfseries #1}, fonttitle=\small\bfseries, coltitle=white,
  colbacktitle=warnline}
\newtcolorbox{okbox}[1]{enhanced, colback=okbg, colframe=okline,
  boxrule=0.6pt, arc=1.5mm, left=3mm, right=3mm, top=2mm, bottom=2mm,
  title={\bfseries #1}, fonttitle=\small\bfseries, coltitle=white,
  colbacktitle=okline}
\newtcolorbox{tipbox}[1]{enhanced, colback=tipbg, colframe=tipline,
  boxrule=0.5pt, arc=1.5mm, left=3mm, right=3mm, top=1.5mm, bottom=1.5mm,
  title={\bfseries #1}, fonttitle=\small\bfseries, coltitle=white,
  colbacktitle=tipline}

\newcommand{\code}[1]{\texttt{#1}}
\DefineVerbatimEnvironment{Verilog}{Verbatim}{frame=single,fontsize=\small}

\begin{document}

% ==================== 封面 ====================
\begin{titlepage}
\centering
\vspace*{2.2cm}
{\color{deep}\Huge\bfseries 外设驱动、摄像头配置与固化}\\[6pt]
{\color{deep}\Large\bfseries 嵌入式对接教程（二）}\\[14pt]
{\large 2026 FPGA 创新设计赛道（安路）赛题二}\\[6pt]
{\large 嵌入式方向 · 04\_embedded}\\[20pt]

\begin{tcolorbox}[enhanced, colback=codebg, colframe=accent, boxrule=0.8pt,
  arc=2mm, width=12.5cm, left=4mm, right=4mm, top=3mm, bottom=3mm]
\centering
{\small 读者：嵌入式方向队员（零基础，建议先读第一册）\\
目标：把状态送出板子（UART）、把摄像头配起来（I2C）、把 bit 烧进去（QSPI）\\
前置要求：会 Verilog、跑过 ModelSim；了解 TD 软件界面\\
预计学习时间：8～12 天（每天 2～3 小时）}
\end{tcolorbox}

\vspace{8mm}
\begin{tcolorbox}[enhanced, colback=tipbg, colframe=tipline, boxrule=0.5pt,
  arc=2mm, width=12.5cm, left=4mm, right=4mm, top=2mm, bottom=2mm]
\centering
{\small\color{deep} \textbf{这一册解决什么}\\
"从 bit 到能演示"的最后一段路：三个外设通道（UART 遥测、
按键）、一个必需的外设配置（摄像头寄存器表）、
一条可复现的固化流程（TD 工程 + QSPI）。}
\end{tcolorbox}

\vfill
{\small\color{black!50} fpga-ptz-tracker 教程系列 · 2026-09}
\end{titlepage}

\tableofcontents
\newpage

% ============================================================
\part{从零学习教程}

% ============================================================
\section{你在造什么}

\subsection{一句话}

\textbf{给这块板子接上"嘴"（UART 遥测）、"耳朵"（按键）、}
\textbf{"眼睛的说明书"（摄像头配置），并让它一上电自己跑起来（固化）。}

\subsection{三个外设通道}

\begin{Verilog}
UART 遥测:  FPGA --(115200 串口)--> 电脑/上位机 (每帧一包目标状态)
按键:       KEY0 -> 切阈值组   KEY1 -> 跟踪使能开关 (经官方消抖)
I2C:        FPGA --(SCL/SDA, 0x36)--> SC500CS 摄像头 (116 条寄存器)
固化:       TD 综合 -> bit -> BitWizard 写 QSPI Flash -> 上电即运行
\end{Verilog}

\subsection{做完是什么效果}

插上 USB 串口，串口助手每 16.7ms 收到一包 9 字节数据
（\code{A5 01 1C 02 53 01 00 F1 5A} 这样的东西），
解析出"已锁定、坐标 (540, 339)"；按 KEY0 红绿蓝阈值组轮流切换；
拔掉电脑，板子独立上电，15 秒内出图并进入跟踪。

\begin{okbox}{学完本册你能做到}
\begin{itemize}
  \item 手算 UART 波特率分频，说出 8N1 帧的每一位
  \item 设计带定界与校验和的遥测协议，并解释"发送中跳过"策略
  \item 讲清 I2C 的 START/ACK/STOP 与寄存器写事务格式
  \item 读懂 SC500CS 720p60 寄存器表与分辨率变更的连锁清单
  \item 按 runbook 走完 TD 工程、引脚、固化、五步联调检查点
\end{itemize}
\end{okbox}

% ============================================================
\section{阶段 0：UART 串口原理（1～2 天）}

\subsection{一根线的协议}

UART 没有时钟线，靠\textbf{双方约定的波特率}同步。一帧 8N1：

\begin{Verilog}
空闲 = 高电平
起始位: 1 拍低       <- 接收方看到下降沿, 知道"要来了"
数据位: 8 位, LSB 先发
停止位: 1 拍高
(无校验位: N = None)
\end{Verilog}

接收方在起始沿之后，每隔\textbf{一个位时间}采一次样；
采在位的正中间最稳（抗抖动）。

\subsection{波特率与分频系数}

FPGA 没有"波特率发生器"，只有系统时钟，用分频凑：

\begin{Verilog}
位时间 = CLK_DIV / f_clk
CLK_DIV = 646, f_clk = 74.25MHz
实际波特率 = 74.25e6 / 646 = 114,938   (目标 115200, 误差 -0.23%)
\end{Verilog}

\begin{keybox}{误差多大算安全}
UART 容忍总误差 ±2～3\%（收发各占一半）。0.23\% 有十倍余量，
随便换时钟都能接受。校验方法：\code{实际值 / 目标值 - 1}，
绝对值小于 2\% 就上。
\end{keybox}

\subsection{练习}

\begin{enumerate}
  \item 手算：50MHz 时钟下，115200 波特率的分频系数是多少？误差多少？
  \item 画出一帧字节 \code{0xA5} 在线上（txd）的电平序列（含起止位）
  \item 思考：如果时钟有 5\% 误差、连发 9 个字节，为什么后几个字节更容易错？
\end{enumerate}

% ============================================================
\section{阶段 1：遥测协议设计（1～2 天）}

\subsection{9 字节一包}

\begin{Verilog}
偏移  内容            说明
0     A5              包头(定界)
1     valid           目标有效位
2     cx 低 8 位
3     cx 高 3 位      (11bit 坐标拆成两个字节)
4     cy 低 8 位
5     cy 高 3 位
6     preset          当前阈值组号
7     checksum        前 7 字节低 8 位累加和
8     5A              包尾(定界)
\end{Verilog}

三个设计点：\textbf{双字节定界}（A5 开头 5A 结尾，上位机扫描同步头）；
\textbf{累加和校验}（一行代码，挡住所有单字节错误）；
\textbf{坐标拆高低字节}（11bit 装不进 1 字节，低 8 + 高 3 是标准做法）。

\subsection{"快照 + 发送中跳过"策略}

每帧（16.7ms）触发一次组帧，但一包要发 780us——

\begin{Verilog}
帧拍到达 && 不在发送中 -> 快照当前坐标, 开始发
帧拍到达 && 正在发送   -> 跳过本次 (下一帧再发)
发送期间一切不变 -> 绝不打断当前包
\end{Verilog}

\begin{warnbox}{为什么不做发送 FIFO}
"跳过"策略保证任意时刻板子只发\textbf{完整包}，
上位机不需要处理残包。代价是发送忙时丢帧——但 780us 远小于
16.7ms 帧周期（占用率约 5\%），实际几乎不丢。
\textbf{能用策略解决的，不要用缓冲解决}——少一个 FIFO 少一份面积和验证。
\end{warnbox}

\subsection{练习}

\begin{enumerate}
  \item 手算：9 字节 x 10 位 x 8.68us/位 约等于多少？一帧发送占用率是多少？
  \item 写一个 Python 上位机解析伪代码（找 A5、收 9 字节、验证 5A 与校验和）
  \item 思考：如果想加一个"帧号"字段，应该插在哪两位之间？为什么？
\end{enumerate}

% ============================================================
\section{阶段 2：I2C 与摄像头配置（2～3 天）}

\subsection{两条线的协议}

I2C 只有 SCL（时钟）和 SDA（数据），且都是\textbf{开漏}：
谁都不会主动拉高，只能拉低或放手；"1"由\textbf{上拉电阻}提供。

\begin{Verilog}
START:  SCL 高电平期间, SDA 由高变低
STOP:   SCL 高电平期间, SDA 由低变高
数据:   SCL 低电平期间变化, 高电平期间稳定 (接收方在上升沿采)
ACK:    第 9 个时钟, 从机把 SDA 拉低表示"收到"
\end{Verilog}

\begin{warnbox}{整条 I2C 最容易错的地方：位相}
"SCL 高时 SDA 下降"既是 START 的定义，也是\textbf{数据位变化}的
违规情形。从机状态机必须把"数据位"和"起始/停止"分清，
否则会出现"地址识别正确、ACK 却总被判成 START"的诡异现象
（本项目真实遇到过，见 Part II 遗留一节）。
\end{warnbox}

\subsection{摄像头寄存器表}

SC500CS 的配置就是一张\textbf{寄存器写入清单}：
720p60 配置共 \textbf{116 条}，写事务格式：

\begin{Verilog}
{0x6C, reg_hi, reg_lo, data}      <- 4 字节一事务
0x6C = 7bit 地址 0x36 左移一位 + 写方向位 0
(地址 0x36 是 SC500 系列出厂默认)
\end{Verilog}

表存在 ROM 里（\code{uics500reg\_720p60.v}，官方提供），
配置引擎（\code{uicfgcs500\_720p} + \code{uii2c}）逐条读出并发送。

\begin{keybox}{换分辨率不是"换一张表"那么简单}
官方 ROM 文件头写着一份连锁清单：MIPI 接收参数、图像宽高、
ISP、DDR 帧缓存、HDMI 时序、HDMI TX 参数、时钟与时序约束——
\textbf{全部要用同一组分辨率帧率}，否则黑屏/花屏/错位。
本项目锁定 720p60 不动，改分辨率要走接口变更流程。
\end{keybox}

\subsection{练习}

\begin{enumerate}
  \item 手算：7bit 地址 0x36 左移一位，加写位（0）后是多少？
  \item 画出写事务 \code{6C 03 5A 12} 在 SDA/SCL 上的完整波形
    （含两次 ACK）
  \item 思考：为什么 I2C 用开漏而不是推挽？（提示：多设备并联）
\end{enumerate}

% ============================================================
\section{阶段 3：按键消抖（1 天）}

\subsection{机械抖动的物理现实}

按键内部是金属弹片，按下/松开瞬间会\textbf{弹跳几毫秒}——
电平在 0/1 之间抖十几下。直接计数会一次按下多切好几档。

\begin{Verilog}
官方消抖: key_remove_shakes
    I_key_in (低有效取反后进) -> 内部连续采样, 稳定若干 ms 才认
    O_key_trig_out: 一次按键 = 一个脉冲 (不是电平!)
用法: if (key0_trig) preset_idx <= preset_idx + 1;   // 按一下加一
\end{Verilog}

\begin{keybox}{触发脉冲 vs 电平的分工}
消抖模块输出的是\textbf{单拍脉冲}（"按下"这个事件），
不是"按住"电平。所以顶层用 if 直接触发动作，
不需要自己再写边沿检测——\textbf{把边沿检测封在消抖模块里}，
顶层逻辑只剩一行。
\end{keybox}

\subsection{练习}

\begin{enumerate}
  \item 估算：74.25MHz 下消抖 10ms 需要多少拍？
  \item 思考：如果按键接的是"按住就连续加"的旋钮场景，脉冲输出还适用吗？
\end{enumerate}

% ============================================================
\section{阶段 4：TD 工具链与固化（2～3 天）}

\subsection{工程组织与文件清单纪律}

\begin{Verilog}
td_prj/ 工程主文件 (prj.al)
src/    全部 RTL (按 scripts/filelist.f 加入)
filelist.f 三段式:
  (1) 视觉方向新增 (13 个 rtl 文件)
  (2) 嵌入式新增 (system_top / uart_telemetry)
  (3) 仿真专用 -- 禁止加入 TD 工程 (替身/TB/sim_video_src 等)
\end{Verilog}

\begin{warnbox}{filelist 的"第三段"是保命的}
\code{divider\_gate\_sim.v}、\code{ph1p\_logic\_eram\_stub.v} 这些
替身文件如果被顺手加进 TD 工程，会和真网表/真原语\textbf{抢同一个
模块名}——综合结果随机、上板行为不可预期。
把"禁止加入"写进清单文件本身，让它在最显眼的地方。
\end{warnbox}

\subsection{TCL 全流程骨架}

\begin{Verilog}
prj_create -name vision_ptz -device PH1P35MDG324 -force
prj_add_files -file_list scripts/filelist.f
prj_set_constraint -sdc ... -adc ...
compile -impl syn_1     ;# 综合
map    -impl syn_1     ;# 映射
pnr    -impl syn_1     ;# 布局布线
bitgen -impl syn_1     ;# 出 .bit
rep_timing -impl syn_1 ;# 看 WNS >= 0
rep_resource -impl syn_1
\end{Verilog}

\begin{warnbox}{TCL 命令名要在 TD 里逐条核对}
上面的骨架按 TD 6.2.1 的习惯编写，各版本命令名/参数有差异——
\textbf{首跑时打开 TD 的 Tcl 控制台，用实际补全逐个核对}。
骨架的价值是流程顺序（综合-映射-布线-bit 的依赖关系），
不是命令拼写。
\end{warnbox}

\subsection{固化 QSPI：让它"上电即运行"}

\begin{Verilog}
1. TD 工具 -> BitWizard -> 目标器件选板载 QSPI Flash
2. 加载 syn_1/vision_ptz.bit -> Boot Image Tool 转换/烧写
3. 烧写完成 -> 断电 -> 拔 USB -> 重新上电
4. 15 秒内应出图并进入跟踪 (参数全部固化在 bit 里)
5. 验收"拔 PC 运行": 显示器独立供电, 无 PC 参与
\end{Verilog}

\begin{keybox}{为什么必须固化（赛事硬约束）}
赛道要求"冷启动直接跟踪"：评委上电就是演示状态，
不允许插电脑加载。这也反过来说明为什么视觉参数
（阈值/增益/限幅）全部做成 RTL parameter——
\textbf{参数进 bit，上电即生效}。
\end{keybox}

\subsection{练习}

\begin{enumerate}
  \item 打开 TD，用 Tcl 控制台验证骨架里的 \code{compile} 命令实际签名
  \item 列出要把 O\_pwm\_x/y 和按键加到 pin 约束时需要的三样信息
    （脚位、电平标准、原理图页码）
  \item 思考：Boot Image Tool 转换前后文件大小变化的原因
\end{enumerate}

% ============================================================
\section{阶段 5：跑通外设测试（1 天）}

\subsection{两个测试}

\begin{Verbatim}
--- tb_uart_telemetry ...
    注入 3 帧 (valid=1/0/1, 坐标各异) -> 解码串行字节流
    # FRAME OK valid=1 cx=540 cy=339 sum OK
    ==> PASS        <- 全部帧正确且含 >=2 帧 LOCK

--- tb_i2c_config ...
    uicfgcs500_720p + uii2c + ROM 对着 I2C BFM 跑
    # CFG cmp 116/116 OK
    ==> 部分通过 (ACK 位相遗留, 见 Part II)
\end{Verbatim}

\begin{okbox}{本册学习清单（全部打勾就过关）}
\begin{itemize}
  \item 能手算 UART 分频与误差，画出 8N1 波形
  \item 能解释遥测协议的定界/校验/拆字节三要素
  \item 能画出 I2C 写事务波形并说明 ACK 位相陷阱
  \item 能说清 filelist 第三段的作用与分辨率变更的连锁清单
  \item 按 runbook 完成一次固化与五步联调检查点
  \item 跑通 tb\_uart\_telemetry，说清 tb\_i2c\_config 遗留的根因
\end{itemize}
\end{okbox}

% ============================================================
% 嵌入式对接教程（二）· Part II：工程实现详解
% 由 emb2_tutorial.tex \input 引入（共用同一份导言区与主题色）
% ============================================================

\part{工程实现详解}

\begin{tipbox}{这一部分写什么}
外设驱动与配置链逐模块讲透：uart\_telemetry 的 bit\_idx 状态机与
v1 重写、uart\_rx\_mon 的"自由运转"监视器、i2c\_bfm\_sc500 的开漏
建模与 ACK 位相遗留、配置链 TB 的对录方法，以及固化 runbook
逐节的执行细节。本册对应迭代记录的第 4、5、6 轮。
\end{tipbox}

% ============================================================
\section{工程总览}

\subsection{外设层文件清单}

\begin{center}
\begin{tabular}{llr l}
\toprule
\rowcolor{rowtint}\textbf{文件} & \textbf{类型} & \textbf{行数} & \textbf{职责} \\
\midrule
\code{rtl/uart\_telemetry.v} & RTL & 94 & 目标状态遥测（115200-8N1） \\
\code{tb/uart\_rx\_mon.v} & 仿真 & 64 & UART 接收监视器（两处复用） \\
\code{tb/tb\_uart\_telemetry.v} & 测试 & 88 & 3 帧注入逐字节验证 \\
\code{rtl/i2c\_bfm\_sc500.v} & 仿真 & 88 & SC500CS I2C 从机 BFM \\
\code{tb/tb\_i2c\_config.v} & 测试 & 93 & 配置链对录（116 条） \\
\midrule
\code{sim/uics500reg\_720p60.v} & 官方 & --- & 116 条寄存器 ROM \\
\code{sim/uii2c.v / uicfgcs500\_720p.v} & 官方 & --- & I2C 控制器与配置引擎 \\
\midrule
\code{scripts/filelist.f} & 脚本 & --- & TD 源文件清单（三段式） \\
\code{scripts/td\_build\_and\_flash.md} & 文档 & --- & 固化 runbook \\
\bottomrule
\end{tabular}
\end{center}

\subsection{外设层的验证策略}

\begin{center}
\begin{tabular}{lll}
\toprule
\rowcolor{rowtint}\textbf{通道} & \textbf{验证方法} & \textbf{状态} \\
\midrule
UART 遥测 & 自收自解（写\code{->}线上\code{->}监视器解码）逐字节比对 & PASS \\
摄像头配置 & 官方配置链 vs BFM 记录，逐事务对录 ROM & 116/116 一致，位相遗留 \\
按键 & 顶层 glue 一行逻辑，随整链仿真覆盖 & 已覆盖 \\
\bottomrule
\end{tabular}
\end{center}

\begin{keybox}{为什么 UART 可以"自收自解"而 I2C 不行}
UART 是\textbf{点对点}：自己发的自己收，闭环即可证明字节正确——
监视器是"独立的第二双眼睛"（独立实现，不共享发送端代码路径）。
I2C 是\textbf{总线}：主控行为（位相/仲裁/时钟拉伸）需要有从机应答
才能走完流程，所以必须写一个 BFM 当对手方——而 BFM 自身的正确性
又需要真机波形佐证（这正是遗留项的由来）。
\end{keybox}

% ============================================================
\section{模块详解：uart\_telemetry}

\subsection{bit\_idx 状态机（v2，重写版）}

发送状态机把"位时间"当节拍：

\begin{Verilog}
if (div_cnt == CLK_DIV - 1) begin          // 一个位时间到
    div_cnt <= 0;
    case (bit_idx)
        4'd0: begin  O_txd <= 1'b0;                 // 起始位
                     shifter <= payload[byte_idx]; end // 装载
        4'd9: begin  O_txd <= 1'b1;                 // 停止位
                     if (byte_idx == NBYTES-1) sending <= 1'b0;
                     else byte_idx <= byte_idx + 1; end
        default: begin O_txd <= shifter[0];          // 数据位 LSB 先行
                       shifter <= {1'b0, shifter[7:1]}; end
    endcase
    bit_idx <= (bit_idx == 4'd9) ? 4'd0 : bit_idx + 4'd1;
end else div_cnt <= div_cnt + 10'd1;
\end{Verilog}

\begin{warnbox}{v1 的失败：装载与移位双写}
第一版把"装载 payload"和"移位输出"写在同一拍的两个分支里，
装载时移位的状态互相覆盖，波形全乱。重写原则：
\textbf{一个节拍只干一件事}——bit\_idx=0 装载（不发数据），
1..8 只移位，9 只出停止位。状态机的每个状态职责单一，
就没有"两件事打架"的空间。
\end{warnbox}

\subsection{组帧快照与跳过}

\begin{Verilog}
tick_d <= I_frame_tick;                        // 边沿检测
if (I_frame_tick && !tick_d && !sending) begin
    payload[0] <= 8'hA5;
    payload[1] <= {7'd0, I_target_valid};
    payload[2] <= I_cx[7:0];   payload[3] <= {5'd0, I_cx[10:8]};
    payload[4] <= I_cy[7:0];   payload[5] <= {5'd0, I_cy[10:8]};
    payload[6] <= {5'd0, I_preset};
    payload[7] <= 8'hA5 + ...（前 6 字节累加和）;
    payload[8] <= 8'h5A;
    sending <= 1'b1;
end
\end{Verilog}

\code{tick\_d} 边沿检测保证"每帧只触发一次"（即使 frame\_tick
因跨域被拉宽）；\code{!sending} 门控实现跳过策略。
校验和在同一拍用组合加法算出来——这拍只有 7 个 8bit 加法，
摆在帧边界上，时序零压力。

\subsection{波特率分频与误差预算}

\begin{Verilog}
CLK_DIV = 646                     // 参数, 74.25e6 / 115200 = 644.5
实际 = 74125000..? 按 74.25e6 / 646 = 114,938 bps
误差 = 114938 / 115200 - 1 = -0.23%   (容差 ±2~3%)
\end{Verilog}

注意这里选的是"整数分频取最接近值"，不追求精确整除——
\textbf{波特率不需要精确，只需要在容差内}。若换 50MHz 时钟，
重算 CLK\_DIV 即可（参数化设计，一处改全局生效）。

% ============================================================
\section{模块详解：uart\_rx\_mon（TB 侧监视器）}

\subsection{两个关键设计}

\begin{Verilog}
// 1) 两级同步器: 串行输入跨进本域
sync <= {sync[0], I_rxd};   wire rxd = sync[1];

// 2) 半位时间后采样, 对齐位中心
if (!rxd && !busy) begin busy <= 1; div <= CLK_DIV/2; ... end
// 3) 假起始重同步: 起始位中心仍为高 -> 判定假起始, 立刻退出
\end{Verilog}

"半位对齐"是 UART 接收的标准姿势：检测到起始\textbf{下降沿}后
等半个位时间再开始按整位采样，采样点就落在每个位的\textbf{中心}，
左右各有半位余量——最抗时钟误差和线上毛刺。

\subsection{v1 的失败：内联监视器漏沿}

第一版在 TB 里内联写监视代码：\code{@(negedge txd)} 等待起始沿、
再阻塞式延迟采数据位。结果字节值逐位错乱——

\begin{warnbox}{事件同步 + 阻塞延迟的经典坑}
\code{@(negedge)} 只对"当前时刻之后"的边沿敏感；而"包内"的
下一位起始沿可能正好发生在阻塞延迟的\textbf{窗口内部}，
被错过之后状态机再也对不齐。修复：重写为\textbf{自由运转状态机}
（每拍看电平、按计数器推进），没有任何 \code{@} 事件同步——
\textbf{时钟节拍是唯一的同步源}。这个教训和 Step3 的
tb\_pwm 重写（任务式同步 -> 连续监视器）是同一课。
\end{warnbox}

\begin{tipbox}{一个监视器，两处复用}
\code{uart\_rx\_mon} 既用在 tb\_uart\_telemetry（验证遥测模块本身），
也用在 tb\_system\_integration（验证整链里的遥测通道）。
独立实现 + 两处复用 = 第三次交叉验证（发送端 FSM、
接收监视器、整链调用者三种视角）。
\end{tipbox}

% ============================================================
\section{模块详解：i2c\_bfm\_sc500}

\subsection{开漏总线的建模}

\begin{Verilog}
assign I_sda = pull ? 1'bz : 1'b0;   // 只能拉低或放手
// TB 侧: pullup(cam_sda);           // "1" 由上拉提供
\end{Verilog}

BFM 不主动驱动高电平——"放手"（z）后由 TB 的上拉（1kΩ 行为模型）
拉高。这是\textbf{忠实建模开漏}的唯一正确方式；
若用 \code{1'b1} 直接驱动，主控和从机同拍驱动相反值时会得到 X，
仿真掩盖真实故障。

\subsection{位相处理与记录}

\begin{Verilog}
wire scl_rise = I_scl && !scl_d;      // 上升沿: 采样/ACK
wire scl_fall = !I_scl && scl_d;      // 下降沿: 数据变化/计数
wire start_c  = I_scl &&  sda_d && !I_sda;   // SCL 高, SDA 下降
wire stop_c   = I_scl && !sda_d &&  I_sda;   // SCL 高, SDA 上升
// 第 9 个上升沿: 地址匹配 -> 拉低 ACK;  第 9 个下降沿: 释放
// 地址匹配: 移位寄存器 == 8'h6C (7bit 0x36 左移 + 写位)
\end{Verilog}

记录数组放在模块内部（\code{reg [7:0] O\_bytes [0:MAX\_BYTES-1]}），
TB 用层次化引用读取——\textbf{Verilog-2001 的端口列表不允许数组}，
这是绕开限制的标准做法（和 system\_top 的遥测层次化引用同源）。

% ============================================================
\section{配置链测试与遗留（诚实的记录）}

\subsection{tb\_i2c\_config 的"对录"方法}

\begin{Verilog}
官方配置引擎 uicfgcs500_720p + uii2c  --驱动--> SCL/SDA
                                              |
                              BFM 记录全部写事务字节
                                              v
期望 ROM uics500reg_720p60 --------------------> 逐条比对
判据: cfg_done 拉高 且 前 116 条写事务 (addr,reg_hi,reg_lo,data)
      与 ROM 完全一致 (AE 追加写只计数不校验)
\end{Verilog}

这个方法的价值：即使 ACK 位相有争议，\textbf{主控侧发出的
寄存器序列本身}已经被独立证明正确（从 ROM 侧逐条确认）——
把"哪一部分可信、哪一部分待定"切得干干净净。

\subsection{遗留：ACK 位相未收敛}

\begin{warnbox}{现象与根因 }
464 次地址识别\textbf{全部正确}（BFM 侧地址匹配逻辑工作正常），
但每字节之后状态被\textbf{自身的 START 检测误判清零}——
BFM 在第 9 个 SCL 上升沿才拉低 SDA 做 ACK，而"SCL 高时 SDA 下降"
正是 START 的定义，于是 ACK 被自己的 START 逻辑吃掉。
两轮位相修复未收敛（改早了波形对不上，改晚了被判 START）。\\[4pt]
\textbf{处置}：明确列为 TD 阶段项——上板用逻辑分析仪抓真机
SCL/SDA 波形对照，半天可收敛；寄存器序列本身已由 ROM 侧确认，
\textbf{外设层的"配置正确性"不依赖这个未收敛项}。
\end{warnbox}

\begin{tipbox}{无硬件条件下的正确姿势}
一个 BFM 位相问题卡住两轮后，最优解不是继续盲试，
而是：\textbf{记录根因假设 + 明确所需证据（真机波形）+ 判定
对交付的影响（不影响）}。工程交付允许带"已知且已界定"的遗留项，
不允许带"说不清"的问题。
\end{tipbox}

% ============================================================
\section{固化 runbook 详解}

\subsection{filelist.f 的三段式}

\begin{Verilog}
# (1) 视觉方向新增: 13 个 rtl (ser96to24 ... vision_top, font_rom)
# (2) 嵌入式新增:  uart_telemetry.v, system_top.v
# (3) 仿真专用 -- 禁止加入 TD 工程:
#     *_stub.v / *_sim.v (替身), sim_video_src.v, frame_dump24.v, tb/*
\end{Verilog}

第（1）段还定义了"哪些视觉文件进综合"：三个生成文件
（\code{hsv\_inv\_lut.v}、\code{font\_rom\_8x16.v}、
\code{marquee\_rom\_64.v}）在内，仿真专用模块不在内。

\subsection{引脚约束增量}

\begin{Verilog}
set_pin_assign {O_pwm_x} -loc {XX} -io_type LVCMOS33
set_pin_assign {O_pwm_y} -loc {XX} -io_type LVCMOS33
set_pin_assign {I_btn_n[0]} -loc {XX} -io_type LVCMOS33
set_pin_assign {I_btn_n[1]} -loc {XX} -io_type LVCMOS33
# 约束文件新增线网: O_pwm_x / O_pwm_y / I_btn_n[1:0] / O_uart_txd
# clock 约束照例程 SDC 不动
\end{Verilog}

脚位从 40Pin DC3 的空闲脚里挑（对照原理图），电平统一 3.3V。
\textbf{舵机信号线接 PWM 脚，舵机电源必须独立 5V 供电并共地}。

\subsection{固化五步与冷启动验收}

\begin{center}
\begin{tabular}{ll}
\toprule
\rowcolor{rowtint}\textbf{步骤} & \textbf{要点} \\
\midrule
1. 打开 BitWizard & 目标器件选\textbf{板载 QSPI Flash}（不是 FPGA SRAM） \\
2. 加载 bit 并转换 & Boot Image Tool，生成 Flash 镜像 \\
3. 烧写 & 用下载器写 Flash \\
4. 断电重上电 & 拔 USB，验证"无 PC 参与" \\
5. 15 秒内出图 & 出图 + 进入跟踪（参数已在 bit 里） \\
\bottomrule
\end{tabular}
\end{center}

\begin{keybox}{五步联调节奏（每步过了才下一步）}
（1）例程原样跑通出图 \code{->}（2）挂 detect\_top：画面无变化 +
ChipWatcher 见 t\_valid \code{->}（3）挂 OSD：出现 TRK0 文字与框
\code{->}（4）接舵机独立电源：duty 变化时画面稳定
\code{->}（5）固化 + 冷启动验收。\\[4pt]
这个顺序的设计逻辑是\textbf{从"官方已知好"逐级加入变量}，
每级只引入一个新东西——出错时责任面永远是一步之宽。
\end{keybox}

% ============================================================
\section{风险清单与交接}

\begin{center}
\begin{tabular}{lp{8.4cm}}
\toprule
\rowcolor{rowtint}\textbf{风险} & \textbf{对策} \\
\midrule
舵机供电 & 独立 5V 共地；画面随舵机动作抖动 = 供电不合格（仿真测不出） \\
AE 拉动颜色 & 演示前锁曝光，或用 8 组阈值现场切换（视觉已备） \\
时序未收敛 & 例程已知现象；挂视觉后重跑时序报告（视觉 <3\% LUT，预期无碍） \\
tuser 链 & 上板"有数据无帧同步"先查 demosaic tuser（第一册界碑） \\
ACK 位相 & 逻辑分析仪抓真机波形，半天收敛（本册遗留项） \\
\bottomrule
\end{tabular}
\end{center}

\begin{okbox}{本册交付清单}
\begin{itemize}
  \item \code{uart\_telemetry.v}：遥测模块（PASS）+ 复用监视器
  \item \code{i2c\_bfm\_sc500.v} + 对录式配置链 TB（116/116 序列确认）
  \item \code{filelist.f} 三段式清单与 \code{td\_build\_and\_flash.md} runbook
  \item 五步联调检查点与四条风险对策（交接给上板阶段）
  \item 一个边界清晰的遗留项（ACK 位相）+ 收敛路径
\end{itemize}
\end{okbox}

% ============================================================
\section{两册总览}

\begin{center}
\begin{tabular}{lll}
\toprule
\rowcolor{rowtint}\textbf{交付层} & \textbf{主张} & \textbf{证据} \\
\midrule
契约层 & 视觉兼容 96bit 契约 & 视觉回归 14 项 \\
挂载层 & 全链可共存（elaboration 零错） & system\_top \\
行为层 & 数据流贯通（受 demosaic 限制） & tb\_system\_integration \\
外设层 & 遥测逐字节正确 & tb\_uart\_telemetry \\
配置层 & 寄存器序列 116 条一致 & tb\_i2c\_config（位相待收敛） \\
固化层 & 上电即运行 & runbook + 五步检查点（上板执行） \\
\bottomrule
\end{tabular}
\end{center}

\begin{tipbox}{嵌入式对接的两册到这里}
第一册管"装进去、跑起来"，第二册管"送出去、配起来、烧进去"。
两册的任务边界就是上板阶段的工作清单——带着这两册和
\code{README\_分工与移交.md} 的勾选表上板，每一步都有据可查。
\end{tipbox}


% ============================================================
\part*{结语}
\clearpage
\section*{总结与下一步}
\addcontentsline{toc}{section}{总结与下一步}

\subsection*{本册学到了什么}
\addcontentsline{toc}{subsection}{本册学到了什么}
\begin{itemize}
  \item UART 8N1 原理、波特率分频与误差预算
  \item 遥测协议三要素（双字节定界/累加和/坐标拆字节）与跳过策略
  \item I2C 开漏总线、START/ACK/STOP 与寄存器写事务
  \item SC500CS 720p60 寄存器表（116 条）与分辨率连锁清单
  \item 按键消抖的触发脉冲用法
  \item TD 工程组织、filelist 纪律、TCL 流程与 QSPI 固化 runbook
  \item 五步联调检查点与四条风险提示
\end{itemize}

\subsection*{下一步做什么}
\addcontentsline{toc}{subsection}{下一步做什么}
\begin{itemize}
  \item \textbf{上板首战}：按五步检查点逐级点亮（例程 -> 检测 ->
    OSD -> 舵机 -> 固化）
  \item \textbf{收敛遗留项}：用逻辑分析仪对照 tb\_i2c\_config 的
    ACK 位相，半天可收敛；恢复 demosaic tuser 链的 TD 侧验证
  \item \textbf{联合验收}：按 README\_分工与移交.md 的勾选表，
    每过一项打一个勾，作为过程证据存档
\end{itemize}

\end{document}
